\documentclass[10pt,a4paper]{amsart}
\usepackage{amsmath,amssymb,amsthm,graphicx}
% For commutative diagrams we use Xy-pic
\usepackage[all]{xy}

% We use 2cell for 2-commutative diagrams.
\xyoption{2cell}
% Theorem environments.
%
\theoremstyle{plain}
\newtheorem{theorem}[subsection]{Theorem}
\newtheorem{proposition}[subsection]{Proposition}
\newtheorem{lemma}[subsection]{Lemma}

\theoremstyle{definition}
\newtheorem{definition}[subsection]{Definition}
\newtheorem{example}[subsection]{Example}
\newtheorem{exercise}[subsection]{Exercise}
\newtheorem{situation}[subsection]{Situation}

\theoremstyle{remark}
\newtheorem{remark}[subsection]{Remark}
\newtheorem{remarks}[subsection]{Remarks}

\numberwithin{equation}{subsection}

% Macros
%
\def\lim{\mathop{\mathrm{lim}}\nolimits}
\def\colim{\mathop{\mathrm{colim}}\nolimits}
\def\Spec{\mathop{\mathrm{Spec}}}
\def\Hom{\mathop{\mathrm{Hom}}\nolimits}
\def\Ext{\mathop{\mathrm{Ext}}\nolimits}
\def\SheafHom{\mathop{\mathcal{H}\!\mathit{om}}\nolimits}
\def\SheafExt{\mathop{\mathcal{E}\!\mathit{xt}}\nolimits}
\def\Sch{\mathit{Sch}}
\def\Mor{\mathop{\mathrm{Mor}}\nolimits}
\def\Ob{\mathop{\mathrm{Ob}}\nolimits}
\def\Sh{\mathop{\mathit{Sh}}\nolimits}
\def\NL{\mathop{N\!L}\nolimits}
\def\CH{\mathop{\mathrm{CH}}\nolimits}
\def\proetale{{pro\text{-}\acute{e}tale}}
\def\etale{{\acute{e}tale}}
\def\QCoh{\mathit{QCoh}}
\def\Ker{\mathop{\mathrm{Ker}}}
\def\Im{\mathop{\mathrm{Im}}}
\def\Coker{\mathop{\mathrm{Coker}}}
\def\Coim{\mathop{\mathrm{Coim}}}

% Boxtimes
%
\DeclareMathSymbol{\boxtimes}{\mathbin}{AMSa}{"02}

%
% Macros for moduli stacks/spaces
%
\def\QCohstack{\mathcal{QC}\!\mathit{oh}}
\def\Cohstack{\mathcal{C}\!\mathit{oh}}
\def\Spacesstack{\mathcal{S}\!\mathit{paces}}
\def\Quotfunctor{\mathrm{Quot}}
\def\Hilbfunctor{\mathrm{Hilb}}
\def\Curvesstack{\mathcal{C}\!\mathit{urves}}
\def\Polarizedstack{\mathcal{P}\!\mathit{olarized}}
\def\Complexesstack{\mathcal{C}\!\mathit{omplexes}}
% \Pic is the operator that assigns to X its picard group, usage \Pic(X)
% \Picardstack_{X/B} denotes the Picard stack of X over B
% \Picardfunctor_{X/B} denotes the Picard functor of X over B
\def\Pic{\mathop{\mathrm{Pic}}\nolimits}
\def\Picardstack{\mathcal{P}\!\mathit{ic}}
\def\Picardfunctor{\mathrm{Pic}}
\def\Deformationcategory{\mathcal{D}\!\mathit{ef}}
\title{Presenting the Hopf algebra of G as a monogenic extension of the Hopf algebra of H}
\author{The Stacks Project Authors}
\date{stacks-project:etale-cohomology.tex:local-snapshot-20261011; GNU FDL 1.2 or later}
\begin{document}
\maketitle
\noindent\textit{Editorial note.} All delivered mathematical text is the author's own, copied from the pinned Stacks Project chapter file etale-cohomology.tex; no statement or proof step was rewritten, reordered or removed.\par\medskip
\noindent\textit{Editorial note.} The proof is the author's complete argument. It opens by reducing to the preceding discussion of the same section (the discussion being a standing set-up of the section rather than a labelled statement), and then verifies the substitution explicitly; nothing further is suppressed.\par\medskip
\noindent\textit{Editorial note.} The label lemma-monogenic-one occurs in no other staged Stacks chapter, so no cross-chapter restatement had to be weighed for this card.\par\medskip
\noindent\textit{Editorial note.} The counted spans contain no label reference and no bibliography citation, so the card carries no navigation stub and no bibliography excerpt.\par\medskip
\noindent\textit{Editorial note.} The preamble is the chapter's own native preamble excerpt already used for this source: the Xy-pic setup (native preamble lines 14-18) and the native theorem declarations with the mathematical macros (native preamble lines 154-217).\par\medskip
\noindent\textit{Allowed standard prerequisites (editorial).}\begin{itemize}
\item The running situation of the section: the group scheme G, its subgroup H and the element Phi(s). Same chapter, the preceding discussion in which the polynomial relation Phi(s) = t and the Hopf algebras of G and H are introduced. Supplies the objects and the polynomial relation on which the whole computation rests.
\item Monogenic extension of a Hopf algebra. Same chapter, the notion of a monogenic extension of a Hopf algebra, i.e. one generated by a single element. Supplies the second assertion of the statement, which is read off from the displayed presentation.
\item Localization A[s, 1/(pi s + 1)] and its presentation by a quotient. Same chapter together with the ambient algebra conventions: adjoining one variable and inverting an element, and the identification of the resulting algebra with a quotient by the relation Phi(s) - t. Supplies the isomorphism displayed in the statement and the rewriting of 1/(pi\textasciicircum{}ell t + 1) carried out in the proof.
\end{itemize}\medskip
\begin{lemma}
\label{lemma-monogenic-one}
We have
$$
A[s, \frac{1}{\pi s + 1}] =
\left(A[t, \frac{1}{\pi^\ell t + 1}]\right)[s]/(\Phi(s) - t)
$$
In particular, the Hopf algebra of $G$ is a monogenic extension
of the Hopf algebra of $H$.
\end{lemma}

\begin{proof}
Follows from the discussion above and the shape of $\Phi(s)$.
In particular, note that using $\Phi(s) = t$ the element
$\frac{1}{\pi^\ell t + 1}$ becomes the element
$\frac{1}{(\pi s + 1)^\ell}$.
\end{proof}

\end{document}
